\subsection{Research gap}

This thesis contributes to three recurring limitations. First, in existing literatures, the dock is almost always assumed static, which reduces the control problem to steering into a fixed pose, with the exception of handling dock motion with dock mounted sensors. No work estimates a dock's wave-frequency motion onboard, through vision alone, from a vehicle that is itself moving and has a close loop control on that estimate. Second, vision-based docking studies are usually evaluated in clear tanks and pools, where turbidity and illumination are controlled or only coarsely varied \cite{liu2017vision,ni2025vision}. The quantitative envelopses across turbidity, range and marker size are largely missing, with Cesar et al.'s tank evaluation of dos Santos as an exception. Third, the learning-based systems that report the strongest underwater robustness assume GPU-class onboard compute and has training sets that are poorly matched to low-cost ROV platforms with limited power and compute.

This thesis addresses these gaps in two parts. Part~A quantifies the dynamic envelope, which has an end-to-end docking of a BlueROV2 against a dock swaying at wave periods while mapping to real sea states, with the full dynamic marker perception pipeline within the loop, including the failure mode of introducing velocity estimation derived from a moving camera. Part~B quantifies the sensing envelop with real underwater imagery, providing a pixel budget for ArUco detection, a measured light attenuation coefficient and a marker-sizing rule derived by adding high turbidity to real frames. Together, these two parts bound where a passive-marker docking system of this class can operate, a characterisation the literatures currently lacks.